\section{Introduction}
\label{sec:introduction}

Tool-using LLM agents execute multi-step plans in which user goals, model reasoning, tool invocations, and environment responses form a single causal chain. This architecture creates an attack surface that is qualitatively different from single-turn prompt injection: an adversary who controls any environment artifact---a web page, an email, a file, or a tool response---can inject instructions that the agent consumes mid-execution, and the resulting damage depends on which downstream tool calls the agent subsequently issues. Benchmarks such as InjecAgent and AgentDojo formalized this threat and showed that even frontier models execute injected instructions at non-trivial rates~\cite{zhan2024injecagent,debenedetti2024agentdojo}, and recent large-scale studies report that many proposed defenses either sacrifice utility or remain bypassable under adaptive attacks~\cite{ma2026autodojo,bhagwatkar2025indirect}.

Existing defenses fall into three broad families: prompt-level hardening, filter-based input sanitization, and system-level designs such as information-flow separation~\cite{chen2024struq,chen2024secalign,debenedetti2024agentdojo}. A common limitation, from the perspective of deployed agents, is that all three act \emph{before or during} a single execution and have no memory of how previous executions failed. Yet the forensic material for a defense is already present: every completed execution leaves a trajectory whose structure---which action consumed which tool output, which state fed which decision---encodes exactly how the compromise happened. This observation motivates a different defense point in the lifecycle: \emph{after} a suspicious execution is recorded, distill it into a compact safety experience, and inject that experience as advisory memory the next time the same task is attempted.

Turning recorded trajectories into defenses raises three technical questions. First, which executions should be postmortemed? Human review does not scale, so an automated triage signal is required. Second, within a flagged trajectory, which part matters? A trajectory may contain dozens of benign steps interleaved with a single unsafe call, and an experience synthesized from the wrong span teaches the wrong lesson. Third, how should localized evidence be converted into a form that helps---rather than over-constrains---an agent at execution time?

We answer these questions with \method{}, a three-stage pipeline built on trajectory graphs. Each recorded trajectory is deterministically mapped into a heterogeneous semantic DAG whose nodes carry the verbatim text of user requests, agent reasoning, tool calls, and environment states, and whose typed edges carry LLM-justified semantic relations. A graph anomaly detector (an edge-aware GAT over frozen sentence embeddings) performs cheap triage at the graph level. A node-risk scorer trained under weak multiple-instance supervision localizes the risky span to a handful of high-risk nodes and their connecting subgraph. An LLM postmortem judge---which never observes ground-truth labels and is explicitly instructed to reject benign and false-positive cases---then converts the localized evidence into a structured experience (risk type, warning signals, recommended policy). Finally, a paired replay protocol re-executes the recorded cases under an execution controller that gates every proposed agent step, once with and once without the injected experience, so that the defense effect is measured strictly within-case.

We evaluate this pipeline on \tool{}, which provides 2{,}218 sanitized trajectories and aligned semantic graphs over 2{,}000 executable virtual-tool tasks across four domains, with outcome-based labels that separate attack exposure from compromised outcomes. Our contributions are:
\begin{itemize}
\item \textbf{Detector-driven defense experiences.} A complete, label-free chain from graph anomaly detection through risk localization to LLM-filtered experience synthesis, with prompt-injection defenses and credential redaction at the evidence boundary.
\item \textbf{A paired replay evaluation protocol.} Deterministic recorded-environment replay in which future agent steps are hidden until their execution checkpoint, with independent judge families and an input fingerprint that pins prompts, models, and case sets to each output directory.
\item \textbf{Empirical results on \tool{}.} Under a fixed 7:1.5:1.5 grouped split with 15\% abnormal ratio and five model seeds, we report detection quality (ROC/PR), paired ASR reductions with McNemar tests, ablations isolating the contribution of detector-derived localization, localization quality, and per-case cost.
\end{itemize}

Section~\ref{sec:related} situates our approach against agent-security benchmarks and defenses. Section~\ref{sec:method} details the pipeline. Section~\ref{sec:setup} and Section~\ref{sec:results} present the experimental setup and results, and Section~\ref{sec:discussion} discusses limitations.
